%VAHENDID: must kast, millest on välja toodud sinine ja kollane juhe ning mille sees on kaks dioodi ja takistit ühendatud nii, nagu on näidatud juuresoleval skeemil; patarei; multimeeter; samasugune diood nagu mustas kastis, mis on kokku ühendatud takistiga takistusega ligikaudu $R = \SI{2,43}{k\ohm}$. Takisti ja dioodi vabade otste külge on ühendamise hõlbustamiseks külge joodetud juhtmejupid ning takisti ja dioodi ühenduskoha külge on võimalik kinnitada juhe või krokodill
Leida võimalikult täpselt mustas kastis olevate takistite $R_1$ ja $R_2$
väärtused!

\begin{center}
\includegraphics[width=0.45\linewidth]{2022-v3g-ge1-yl.png}
\end{center}

\textit{Vihje:} dioodi läbiv vool sõltub rakendatud pingest mittelineaarselt,
kusjuures vastupinge rakendamisel (st siis, kui vool peaks liikuma noolele
vastassuunas) on voolutugevus tühiselt väike.

\textbf{NB!} Mitte ühendada dioodi otse patarei klemmidele, selle tulemusel
hävineb diood või põleb sisse patareipesa.
